\section{Configuration-Dependent Buffer Safety}
\label{sec:btd_reference}

\subsection{Declared reference and exact recurrence}
The analysis uses a connected, undirected, unit-weight graph with
Laplacian $L$, common capacity $\beta_{\max}$, zero physical link latency, zero
control-execution delay, fluid
occupancies and synchronous physical sampling $t_k=kH$. Corrections
are zero-order held. Initial phases are equal, and targets satisfy
$\beta^*_{ji}=\lambda_{ji}$ for compatible additive markings. These
are reference assumptions, not properties of actual asynchronous
Callisto. The endpoint recurrence is exact for continuous held flow,
not Euler integration or a discrete firing-count plant.

Set $\delta\omega=\omega^u-\omega_{\rm ref}\mathbf1$,
$x_k=\theta(t_k)-\theta(0)-kH\omega_{\rm ref}\mathbf1$,
$w_k=\delta\omega+c_k$ and $v_k=Hw_{k-1}$ for $k\ge1$.
Then $E_k=-Lx_k$ and $\beta_{ji}(t_k)-\beta^*_{ji}=x_{j,k}-x_{i,k}$.
With $D_{\alpha,k}=\operatorname{diag}(\alpha_{i,k})$, the calibrated
law yields
\begin{align}
 v_{k+1}&=v_k-D_{\alpha,k}(bLx_k+aLv_k),\nonumber\\
 x_{k+1}&=x_k+v_{k+1},\quad a=Hk_p,\quad b=Hk_ih.
 \label{eq:btd_full_recurrence}
\end{align}
The integration scale $h$ and physical period $H$ remain distinct.
Zero controller memory and centered startup give $c_0=0$ and the
genuine bootstrap $x_1=v_1=H\delta\omega$; stored previous error is
not a previous plant interval. Appendix~\ref{app:theory} derives these relations.

\subsection{Invariant domain and safety certificate}
Write $L=U\Lambda U^\top$ with $U^\top\mathbf1=0$ and
$U^\top U=I$. The positive diagonal matrix $\Lambda$ contains
the nonzero graph eigenvalues $\mu_2,\ldots,\mu_n$.
Define $\xi=\Lambda^{1/2}U^\top x$, $\eta=\Lambda^{1/2}U^\top v$
and $y=(\xi^\top,\eta^\top)^\top$. The state-dependent symmetric
matrix $M_k=\Lambda^{1/2}U^\top D_{\alpha,k}U\Lambda^{1/2}$ has
eigenvalues in $[\nu_-,\nu_+]=[\underline\alpha\mu_2,
\overline\alpha\mu_n]$. No commutation between $D_\alpha$ and $L$
is assumed. Define
\begin{equation}
 B(\nu)=\begin{bmatrix}1-b\nu&1-a\nu\\-b\nu&1-a\nu\end{bmatrix},
 \quad V_k=y_k^\top(\bar P\otimes I)y_k.
\end{equation}
Here $\bar P\succ0$ is a $2\times2$ common Lyapunov witness and
$\otimes$ denotes the Kronecker product.

\begin{btdlemma}[Endpoint contraction]
\label{lem:btd_endpoint}
If $0<\rho<1$ and
\begin{equation}
 B(\nu_\pm)^\top\bar P B(\nu_\pm)\preceq\rho\bar P,
 \label{eq:btd_d_endpoints}
\end{equation}
then every update in the regional sector satisfies
$V_{k+1}\le\rho V_k$, including time- and node-dependent factors.
\end{btdlemma}
Affine dependence of $B$ gives matrix convexity over the whole
endpoint interval. Orthogonal diagonalization of each symmetric
$M_k$ preserves $\bar P\otimes I$, so one witness covers changing
node factors. Appendix~\ref{app:theory} supplies the full proof.

The importance of a common witness is that the nonlinear surface
changes its factors at every update. Stability of one frozen-gain
characteristic polynomial would not certify the resulting trajectory.
The symmetric projection handles the node dependence without assuming
that all factors or graph modes evolve independently. However, this
certificate remains conditional on the inputs staying in their derived
sector; the following radius test closes that condition.

For $\gamma_i=100/(d_i\beta_{\max})$, define
\begin{equation}
 \Psi_i=\gamma_i^2d_i\max_{\sigma=\pm1}
 [K_e,\sigma K_{\Delta e}]\bar P^{-1}
 [K_e,\sigma K_{\Delta e}]^\top.
\end{equation}
If $V_1\le R$ and $R\Psi_i\le r^2$ for every node, the energy
output bound gives $|X_i|+|Y_i|\le r$. Together with contraction,
induction proves input-domain invariance rather than assuming it.
Let $L^\dagger$ be the Laplacian pseudoinverse and
$R_{ij}^{\rm eff}=(\mathbf e_i-\mathbf e_j)^\top L^\dagger
(\mathbf e_i-\mathbf e_j)$ the effective resistance. Put
$R_{\max}^{\rm eff}=\max_{i\ne j}R_{ij}^{\rm eff}$ and
$m_e=\min(\beta_e^*,\beta_{\max}-\beta_e^*)$. Then
\begin{align}
 |\beta_e-\beta_e^*|^2&\le R[\bar P^{-1}]_{11}R_e^{\rm eff},\nonumber\\
 \Delta_\omega(t)&\le D(R)\rho^{k/2},\quad t\in[t_k,t_{k+1}),\nonumber\\
 D(R)&=H^{-1}\sqrt{R[\bar P^{-1}]_{22}R_{\max}^{\rm eff}}.
 \label{eq:btd_output_bounds_short}
\end{align}
Here $R_e^{\rm eff}$ uses the endpoints of edge $e$.
Affine intersample occupancy covers all physical times.

The unprojected common frequency also needs a bound. Set
$g=(\kappa_I,k_p)^\top$,
$s_\alpha=(\overline\alpha-\underline\alpha)/2$,
$\bar w_0=n^{-1}\mathbf1^\top\delta\omega$ and
\begin{align}
 W_c&=|\bar w_0|+\frac{2s_\alpha}{1-\rho}
       \sqrt{\frac{\mu_nR}{n}g^\top\bar P^{-1}g},\nonumber\\
 W_d&=H^{-1}\max_i\sqrt{R[\bar P^{-1}]_{22}(L^\dagger)_{ii}}.
\end{align}
The following joint sufficient conditions close capacity, limiter
and positive-frequency requirements:
\begin{align}
 &V_1\le R,\quad R\Psi_i\le r^2,\quad 0<\beta_e^*<\beta_{\max},\nonumber\\
 &R[\bar P^{-1}]_{11}R_e^{\rm eff}<m_e^2,\nonumber\\
 &W_c+W_d+\|\delta\omega\|_\infty<c_{\max},\nonumber\\
 &\omega_{\rm ref}-W_c-W_d>0.
 \label{eq:btd_joint_sufficient_conditions}
\end{align}
\begin{btdcorollary}[Reference safety and convergence]
\label{cor:btd_buffer_safety}
Under the reference assumptions, endpoint certificate and joint
conditions, all buffers remain strictly within capacity, the limiter
is inactive and physical frequencies stay positive. Buffer deviations
and pairwise frequency differences converge to zero. If a separate
tolerance $\Delta_{\rm safe}$ is specified, $D(R)<\Delta_{\rm safe}$
also certifies its all-time bound.
\end{btdcorollary}
Summable common-frequency increments give a finite limiting rate,
which need not equal the mean uncontrolled frequency for Fuzzy. A failed sufficient
test is inconclusive; these conditions are not necessary stability
boundaries or a certificate for delayed integer execution.

\subsection{How the parameter constraints interact}
The certificate exposes relationships between initial mismatch,
buffer reserve, graph structure and controller scales, rather than
independent recommended ranges for each parameter. Define
$c_P=[1,1]\bar P[1,1]^\top$. At the centered bootstrap,
\begin{equation}
 V_1=H^2c_P\,\delta\omega^\top L\delta\omega.
\end{equation}
Thus physical observation time and the mismatch direction enter the
same initial-energy test. A statement about one maximum frequency
difference alone does not identify this energy on an arbitrary graph.
In the worked triangle $c_P=1$, but it need not equal one for another
witness. The joint bounds must then accommodate both this initial
energy and the nonlinear input region.

\begin{table*}[t]
\caption{Coupled roles of system parameters in the declared reference certificate.
Every row must be evaluated for the same controller, graph and initialization;
the table is not a set of independent necessary stability limits.}
\label{tab:parameter_roles}
\centering
\begin{tabular}{@{}lp{0.72\textwidth}@{}}\toprule
Parameter group & Role in the sufficient conditions\\\midrule
Initial clock frequency mismatch & $H^2c_P\delta\omega^\top L\delta\omega\le R$ sets the required radius.\\
Initial occupancy and capacity & Compatible target $0<\beta_e^*<\beta_{\max}$ and reserve $m_e$ must contain the output bound.\\
Neighbour count and normalization & $\gamma_i=100/(d_i\beta_{\max})$ enters $\Psi_i$, together with graph degree and Fuzzy input scales.\\
Topology & $\mu_2,\mu_n$ determine spectral endpoints; effective resistance determines edge and pair outputs.\\
Control timing and gains & $a=Hk_p$, $b=Hk_ih$ enter the common endpoint inequalities; changing $H$ changes the certificate.\\
Correction and absolute frequency & $W_c+W_d+\|\delta\omega\|_\infty<c_{\max}$ and $\omega_{\rm ref}>W_c+W_d$ close execution constraints.\\
Marking and path history & Targets must match initialization; delayed in-flight frames cannot be silently removed when forming the zero-latency reference.\\\bottomrule
\end{tabular}
\end{table*}

For example, increasing capacity at unchanged Fuzzy gains changes
$\gamma_i$ and therefore the frame-count input scaling. Changing
the reference sampling interval changes $a$ and $b$; reconfiguring
local polling can additionally change $h$. These distinct effects
must not be replaced by a change to one nominal gain alone. Adding neighbors
changes the measured sum, normalization and graph spectrum together.
Table~\ref{tab:parameter_roles} explains why consistent parameter
changes require recertification, even when one individual inequality
appears less restrictive.

The initial occupancy assumption is equally important. This theorem
starts at compatible targets with zero controller memory. An arbitrary
non-centered start introduces an affine error or different phase
coordinates and can alter the first controller call. It cannot be
certified by inserting a new occupancy value into the same centered
bootstrap. Likewise, finite link latency/control-execution delay and asynchronous
polling belong to the Callisto simulations here, not to a hidden
delay margin in the zero-latency proof. These distinctions allow the
mathematical conditions to be used as reproducible design tests
without presenting them as absolute hardware limits.

\subsection{Strict headroom comparison}
For the triangle use $\beta_{\max}=2048$, compatible centered marking 1024,
$\omega^u=(1.0002,0.9998,0.9999)^\top$,
$\omega_{\rm ref}=29999/30000$, $H=h=30000000/29999$ and
$(k_p,k_i)=(2\times10^{-6},5\times10^{-11})$.
The matched Fuzzy gains are
$(K_e,K_{\Delta e},K_u)=(0.1,29999/7500,-768/18749375)$.
A witness is
\begin{equation}
 \bar P=\frac1{50000}\begin{bmatrix}9&181\\181&49629\end{bmatrix},
 \quad\rho=1599/1600,\quad r=1.
\end{equation}
Exact rational endpoint checks and outward enclosures verify the
joint conditions. A coarse excursion bound is 32.237442 frames;
sharper interval propagation and a contracting infinite tail give:
\begin{btdtheorem}[Prescribed-initialization headroom benefit]
\label{thm:btd_headroom}
For this reference and zero controller constructors, both methods
satisfy the safety and convergence conclusions above. Their
all-time excursions $B_J=\sup_{t\ge0,e}|\beta_e^J(t)-1024|$ obey
\begin{equation}
 B_{\rm F}<17.315<23.40370130812701\le B_{\rm PI}.
 \label{eq:btd_strict_comparison}
\end{equation}
Minimum two-sided Fuzzy headroom exceeds 1006.685 frames and
improves on PI by more than 6.08870130812701 frames.
\end{btdtheorem}
\begin{proof}[Proof outline]
PI reduces to an exact rational scalar recurrence with an attained
prefix peak at sample 111. Outward interval propagation encloses the
full nonlinear Fuzzy prefix through sample 200. Tail bounds of
17.958057 frames (PI) and 12.292711 frames (Fuzzy) are below their
respective prefix peaks. Affine intersample flow excludes a missed
larger peak. Appendix~\ref{app:theory} contains the rational witness,
interval arithmetic and complete proof.
\end{proof}
\paragraph{What closes the infinite horizon.}
The initial energy is the exact fraction
$V_1=234000000/899940001$. With this radius the analyzed inputs
satisfy $|X|+|Y|<0.171455<1$, requested correction remains below
0.004166148 and physical frequency stays above 0.99603385.
These bounds establish that the unconstrained trajectory never reaches
the limiter, closing an assumption that otherwise could be circular.
The common-frequency bound is also necessary: a network of equally
nonpositive rates would have zero frequency spread but invalid progress.

The strict ordering compares Fuzzy's certified upper bound with PI's
attained prefix peak. Both infinite-tail bounds are below their
prefix peaks, and intersample occupancy is affine. Comparing only
conservative upper bounds would not establish this ordering;
Appendix~\ref{app:theory} provides the complete arithmetic.

\begin{table}[t]
\caption{Certified prefix and tail roles in the strict reference comparison.
All entries are fluid frames; prefix enclosures are outward-rounded here.}
\label{tab:btd_prefix_tail}
\centering
\begin{tabular}{@{}lll@{}}\toprule
Controller & Prefix peak enclosure & Tail upper bound\\\midrule
PI & $[23.40370130812701,$ & $17.95805698665804$\\
 & $\phantom{[}23.40370130812702]$ & \\
Fuzzy & $[17.314056817962,$ & $12.29271005372874$\\
 & $\phantom{[}17.314056817963]$ & \\\bottomrule
\end{tabular}
\end{table}

The excursion theorem fixes one initial frequency vector and matched
gains. The following argument establishes a safety advantage for a
family of configurations at fixed capacity and normalization.
It does not compare independently optimized controllers.

% Staged proof content; integrated only after the simulation handoff is checked.
\subsection{Initial mismatch and the admissible occupancy interval}
\label{subsec:btd_occupancy_domain}
The joint conditions define a sufficient domain of configurations.
For the centered bootstrap, let
$Q=\delta\omega^\top L\delta\omega$ and
$S_x=[\bar P^{-1}]_{11}$. Since each undirected connection is
counted once in $L$,
\begin{equation}
 Q=\sum_{\{i,j\}\in\mathcal E_u}(\omega_i^u-\omega_j^u)^2,
 \qquad V_1=H^2c_PQ.
 \label{eq:btd_mismatch_energy}
\end{equation}
Here $\mathcal E_u$ is the undirected edge set. Choosing $R=V_1$
makes the buffer and input-region tests explicit:
\begin{equation}
 H^2c_PQ S_xR_e^{\rm eff}<m_e^2,\qquad
 H^2c_PQ\Psi_i\le r^2.
 \label{eq:btd_configuration_slice}
\end{equation}
These tests require the endpoint, correction and positive-frequency
conditions above. They show that initial mismatch and available
storage are coupled; neither a frequency range nor a buffer capacity
alone specifies a safe configuration. In particular, equal maximum
initial frequency spreads can have different $Q$.

A sharper configuration comparison follows from the certified
excursions. Use $z$ for a common initial occupancy and target,
distinguishing it from the recurrence coefficient $b=Hk_ih$.
Keep capacity, gains, frequency vector, initial phases and timing
fixed, and set all compatible zero-delay markings to $z$.

\begin{btdcorollary}[A strict advantage in admissible occupancy]
\label{cor:btd_occupancy_advantage}
For the prescribed reference comparison, changing the common initial
occupancy, target and compatible marking from 1024 to $z$ leaves
each controller's phase and frequency trajectories unchanged.
The interval of targets with uniformly positive two-sided headroom is
\begin{equation}
 \mathcal I_J=(B_J,\beta_{\max}-B_J),\quad J\in\{\mathrm{PI},\mathrm F\}.
 \label{eq:btd_target_interval}
\end{equation}
Consequently $\mathcal I_{\rm PI}\subsetneq\mathcal I_{\rm F}$.
At $\beta_{\max}=2048$ and $z=20$, Fuzzy has all-time occupancy in
$(2.685,37.315)$, whereas the unprotected PI reference crosses
the lower capacity boundary.
\end{btdcorollary}
\begin{proof}
In the unprotected fluid reference,
$\beta_{ji}^{J,z}=z+x_j^J-x_i^J$.
Subtracting target $z$ cancels the shift from the measured errors.
Fixed $\beta_{\max}$ preserves $\gamma_i$, so induction over controller updates
preserves the inputs, memory and phase trajectories. Reverse channels
have opposite deviations; the minimum two-sided headroom is therefore
$\min(z,\beta_{\max}-z)-B_J$. The strict excursion ordering proves the interval
inclusion and the stated safe/crossing case. The precision of the
prefix/tail certificate, rather than the coarse Lyapunov buffer
bound, establishes the Fuzzy capacity conclusion here.
\end{proof}

The same example connects the margin benefit to tolerable initial
mismatch. Along $\delta\omega(s)=s(7,-5,-2)^\top/30000$, the
initial spread is $400s$ ppm. In the non-clipped PI reference,
linearity gives $B_{\rm PI}(s)=sB_{\rm PI}(1)$. For $z=20$ its
capacity threshold is approximately $341.8263$ ppm, while Fuzzy
is certified safe at the 400-ppm point. This establishes an advantage
along this fixed mismatch direction. It does not certify every
400-ppm vector or a uniform Fuzzy bound for all smaller $s$.
Appendix~\ref{detail:subsec:btd_occupancy_translation} gives the full
argument and parameter-domain interpretation.

The advantage also holds on a quantified parameter interval. Keeping
the mismatch direction fixed, outward interval propagation on all
20 cells covering $s\in[0.999,1.001]$ gives a uniform Fuzzy
all-time excursion below 17.342 frames. PI's attained peak is above
23.3802976 frames throughout this interval. Thus the Fuzzy reference
is safe at $z=20$ for every initial spread from 399.6 to 400.4 ppm,
while each corresponding unprotected PI trajectory crosses zero.
This certifies a region in the $(s,z)$ configuration slice, rather
than inferring an interval from a finite grid of successful runs.



For an integer-valued compatible prefill $z$ in this fluid reference,
the bounds have a direct interpretation. Over the proved parameter
interval, the minimum common prefill with positive uniform reserve
is 18 for Fuzzy and 24 for PI. Fuzzy's bound is below 18, but 17
fails at the included $s=1$ point; PI's exact scaled peak lies
strictly between 23 and 24 throughout the interval. This concerns
initial/target occupancy at fixed capacity, not a reduction in the
configured physical buffer capacity.
